\begin{textsample}{POS Dim 4 – df – Score 8.00 – df/df\_inf\_2.txt}  \label{ex:f4_pos_227}
2.
CONTINUANDO NOSSA CONVERSA : a \textbf{identidade} do Distrito Federal expressa no Currículo Linhas modernas.
Desenho futurístico.
Leveza expressa com concreto armado.
Projeto arquitetônico.
Patrimônio cultural da humanidade.
Brasília, \textbf{capital} do Brasil.
Será só isso?
Falar do Distrito Federal vai muito além de expor Brasília como projeto de cidade planejada que se materializou em realidade no Planalto Central, reconhecida por suas formas singulares de arquitetura e considerada patrimônio da humanidade pela Organização das Nações Unidas para a Educação, a Ciência e a Cultura – UNESCO.
Mas por que discorrer sobre o Distrito Federal?
Porque discutir currículo envolve a reflexão sobre o \textbf{território}.
De que lugar que se fala?
Como ele se constitui histórica e socialmente?
Quais são suas singularidades e como elas se expressam?
Como se revela sua paisagem natural e culturalmente constituída?
Que grupos sociais se encontram nesse \textbf{território}?
Que crianças e quais infâncias pertencem a esse lugar e como elas se relacionam?
Todas essas questões perpassam a constituição de um currículo.
Se tais reflexões fossem ignoradas não haveria a necessidade de cada unidade da federação elaborar seu próprio currículo.
São as peculiaridades de cada \textbf{território} que dão subsídio para a formulação do currículo que contribui para o planejamento, para a prática pedagógica e para a avaliação do processo educativo.
É primordial pensar sobre o \textbf{território} que se constitui, influencia a atuação pedagógica e que, dialeticamente, sofre a influência de todos os educadores.
O Distrito Federal está situado na Região Centro-Oeste, sendo uma das 27 unidades da federação brasileira.
É a menor unidade federativa do Brasil e a única que não possui municípios.
Em seu \textbf{território}, encontra -se a \textbf{capital} federal do Brasil, Brasília, que também é a sede do governo do Distrito Federal.
O Distrito Federal é um quadradinho no meio do mapa do Brasil, mas nem por isso deixa de ter suas características próprias.
É constituído por pessoas de todas as unidades federativas do Brasil, bem como por estrangeiros que, por inúmeros motivos, residem em seu \textbf{território}.
Mesmo sendo o ente federado de menor idade, possui uma \textbf{população} que nasceu, cresceu e contribui para a constituição social, cultural e histórica desse lugar.
As \textbf{populações} \textbf{indígenas} já residem nesse \textbf{território} há muito tempo.
Há quem diga que no Distrito Federal não tem \textbf{indígenas}, sendo esta uma visão naturalizada pelo senso-comum que precisa ser desmistificada para que esses povos sejam respeitados e considerados nas políticas públicas, sobretudo as que versam sobre a educação.
Há estudos que revelam que, no século XVI, o \textbf{território} central do país, que inclui o atual Distrito Federal, era ocupado por \textbf{indígenas} do tronco linguístico Macro-jê, como os Acroás, os Xacriabás, os Xavantes, os Caiapós, os Javaés etc. ( CHAIM, 1983 ) e, atualmente, “ no Distrito Federal residem três etnias, a saber : Tapuyas/Fulni-ô, Tuxa e Cariri Xocó ” ( NASCIMENTO, 2018, p. 22 ).
Além dos povos \textbf{indígenas}, é \textbf{marcante} a presença de comunidade remanescente quilombola, a saber a Comunidade Mesquita, que fica na região do entorno do Distrito Federal.
É importante destacar esse ponto, pois o Distrito Federal tem intensa relação, das mais diversas formas, com as cidades de seu entorno.
Ademais, notam -se comunidades que residem no campo, assentados e acampados da reforma agrária, povos tradicionais, entre outros.
Toda essa diversidade humana presente revela a constituição de seu peculiar \textbf{território}.
Apesar de Brasília ser sonhada como uma cidade planejada, com linhas arquitetônicas desenhadas em concreto armado, ao seu redor, o Distrito Federal cresceu, criando alternativas próprias para a realidade que emergia dos agrupamentos sociais que já existiam aqui, que vieram para cá ou que nasceram aqui.
Nessa expansão territorial, criou -se 31 regiões administrativas ( DISTRITO FEDERAL, 2018c ), cada uma com sua \textbf{identidade}, organização e necessidades próprias.
O Distrito Federal é caracterizado pela beleza natural do Cerrado, sendo este o segundo maior bioma do país, com cerca de um terço da sua biodiversidade.
É marcado principalmente pelo clima tropical, com uma estiagem que se prolonga por aproximadamente cinco meses.
Uma reflexão que o Cerrado provoca, sobretudo quando se pensa uma educação para a sustentabilidade, como prevê um dos Eixos Transversais deste Currículo, é acerca dos impactos causados pela perda da biodiversidade, imposta pelo adensamento populacional e a expansão da agropecuária. > Até a década de 60, a região do Cerrado era considerada como marginal para a agricultura intensiva.
Sua ocupação foi motivada principalmente pelas mudanças na estrutura rodoviária iniciada com a implantação de Brasília e pela criação do Programa de Desenvolvimento do Centro-Oeste ( Polocentro ), na década de 70, que levou a uma intensa \textbf{migração} em busca de terras a custos mais baixos em relação ao \textbf{sul} do país e incentivos fiscais para abertura de novas áreas \textbf{agrícolas}.
Como resultado dessa política, grandes áreas de Cerrado foram desmatadas, sendo o Bioma Cerrado atualmente considerado como um dos 25 ecossistemas do planeta, com alta biodiversidade, que está ameaçado ( DISTRITO FEDERAL, 2018b, s.p. ).
Portanto, são muitas as singularidades que constituem um \textbf{território} e que não cabem nessas poucas páginas, mas que já são capazes de provocar uma reflexão acerca do que constitui esse lugar, seus povos, seu ambiente natural e cultural e sobre como esses elementos precisam ser pensados na vivência de uma prática educativa que possui uma perspectiva de educação integral.

% matched lemmas: agrícola, capital, identidade, indígena, marcante, migração, população, sul, território
\end{textsample}
